\section{What this continuation adds}
\label{sec:scope}

The practical question is how to reduce the cost of verified topology
operations while removing genuine combinatorial search barriers.  The
repository contains many useful exact operators; a new report must identify
which obstacle it actually removes.  This continuation contributes a new
parameterized theorem for choosing dual forests and an independently tested
improvement to interval-certificate replay.  The theorem and implementation
operate at different interfaces, which makes their scope straightforward to
audit.

\subsection{The inspected checkpoint}

We inspected both requested directories at the pinned revision, including
all five ZIP archives then present in \code{docs/incoming}.  These contain
component-gauge cut atlases, dual-forest normalization, dual forests with
homological progress, flag regularization, and positive cancellation.
They are listed with hashes in the package provenance file.  Their results
are antecedents, not new results of this continuation.  In particular the
following ingredients were already available:

\begin{enumerate}
\item exact column-leaf peeling and its greatest stopping-set obstruction;
\item source-authenticated removal of native three-handles by a rooted
  dual forest;
\item an exact tree/stopping-set criterion, and a polynomial
  graphic--partition matroid-intersection algorithm when each useful
  column has at most two incident rows;
\item the fact that an original native cocore has at most three actual
  occurrences, isolating $(1,1,1)$ as the remaining native support type;
\item saturated interval-bundle compilation and independently replayable
  whole-cut handlebody certificates.
\end{enumerate}

The maintained whole-cut corpus has 106 cases.  At this pin its existing
producer returns 92 complete handlebody-census certificates and 14
inconclusive results.  The behavior-preserving performance patch preserves
that classification.  The optional lower-forest portfolio adds five
independently verified whole-cut positives.  The separate forest research
operator decides a more
general supplied graph/matrix problem; it is not inserted into the native
recognition schedule without a source adapter and a justified dispatch rule.
The detailed measurements and fresh source checks appear in
\cref{sec:experiments}.

\subsection{A new parameter, followed by a stronger algorithm}

The preceding two-row theorem is polynomial even when the dual graph has
many vertices.  Its limitation is the presence of columns that encode a
genuine conjunction: two other rows must disappear before the third can
peel.  Treating either premise as sufficient would be unsound.

Our first extension guesses the exceptional activations and reduces their
prerequisites to a bounded number of overlapping edge-cover requirements.
This gives an exact deterministic algorithm, with the exceptional-column
count in the exponent instead of the dual rank.  It also gives a useful
reference implementation and a transparent completeness proof.

A stronger reduction removes the need to guess an activation order.  Only
the rows occurring in nonlinear rules need to remember the influence of
the tree.  Everything else is ordinary directed reachability.  The effect
of a spanning tree on these exceptional rows is a finite bit mask, and
successful masks form a monotone Boolean function that can be evaluated
by forward chaining.  We compile that function into a determinant filter.
Unlike an inclusion--exclusion sum that asks only whether every color
appears, the filter accepts precisely the masks that permit complete
nonlinear closure.  Each valid copied common-basis monomial retains its
original nonzero determinant coefficient: the filter multiplies it by one.
Invalid monomials are multiplied by zero.

\begin{center}
\begin{tabular}{P{0.26\textwidth}P{0.28\textwidth}P{0.36\textwidth}}
\toprule
Result & Time, ignoring polynomial factors & Exact scope\\
\midrule
Preceding two-row theorem & $1$ & Every useful column has support at most two.\\
Deterministic extension & $3^k k!(m+1)^{2k}$ & At most $k$ exceptional triple-unit columns.\\
Ideal determinant filter & $J\leq2^d\leq8^k$ & $d$ nonfree special SCCs; one-sided randomization.\\
Deterministic width algorithm & $J(m+1)^w$ & $w$ is their implication-poset width.\\
Independent positive check & Polynomial & Supplied tree and actual unsigned peel order.\\
\bottomrule
\end{tabular}
\end{center}

Here $m$ is the number of rows, $k$ counts exceptional columns, $d$ counts
their nonfree ordinary SCCs, and $J$ counts the ideals of the induced
reachability poset.  Every ideal is a downward-closed set.  Verification
does not trust the Boolean table, the determinant evaluation, the random
seed, or the discovery program.  It checks the original finite input directly.

\subsection{Where the implementation gains occur}

The interval-orbit producer already used prefix-gap information in one
weighted contraction path.  Three independent certificate replay paths
still computed the shift of each endpoint by scanning all preceding gaps.
This was a real repeated operation in whole-cut verification.  We supply
a verifier-owned prefix implementation, preserve the exact quotient and
inverse conventions, and compare against the old sums on complete
certificates and independently generated finite examples.

A separate change shares internally reconstructed geometry within one
composed call.  On the producer side, the shared object is its own geometry.
On the verifier side, it is the verifier's fresh reference reconstruction.
No public prepared-object parameter or persistent cache is introduced.
This distinction prevents an apparent optimization from quietly shifting
the trust boundary to producer-supplied data.

\subsection{Relation to Lackenby's announced target}

The primary sources use two displayed quasi-polynomial exponents.  Oxford's
2021 announcement states $n^{c\log n}$, whereas the detailed Oxford handout
displays $n^{O((\log n)^2)}$.  Both are quasi-polynomial; we do not silently
identify them \cite{Oxford2021,LackenbyTalk2021}.  The July 2026 preprint
supplies a hierarchy algorithm, restart analysis, and compressed certification
results; it does not establish the announced quasi-polynomial bound
\cite{Lackenby2026}.

Our main theorem has an especially useful local consequence: if a supplied
state has polynomial encoded size and only $O(\log n)$ exceptional rows,
the randomized forest subproblem has polynomial witness-discovery time.
More generally, bounded implication width gives a deterministic polynomial
local algorithm, and logarithmic width gives a deterministic quasi-polynomial
local algorithm.
Proving that the whole hierarchy always reaches such states, or that
there are only quasi-polynomially many states of controlled size, remains
separate work.  A verified local improvement is a concrete component of
that program, not a substitute for its global argument.

\subsection{Attribution and novelty}

Directed reachability, strongly connected components, matroid intersection,
Cauchy--Binet, Boolean M\"obius inversion, and polynomial identity testing
are established tools.  The standard common-basis polynomial and color
sieve are explicitly treated in the algebraic matroid literature
\cite{Harvey2009,EibenKoanaWahlstrom2025}.  The contribution here is the
exact reduction of forest-compatible peeling to a bounded exceptional-row
state, its Boolean determinant compilation, its integration with the
native three-occurrence theorem, the implication-poset refinement, and
the supplied audited implementations.  Efficient lattice transforms are
also classical; the relevant bound is credited to
Bj\"orklund et al.\ \cite{FastZeta2012}.
We make no exhaustive literature-priority claim.  All claims of progress
are relative to the inspected repository and the cited mathematical inputs.
